\section{Experiments: When Does Escalation Pay?}\label{sec:experiments}

\paragraph{Protocol.}
Rungs: P = the small judge pairwise with confidence $|\lambda|$; S = the same judge pointwise with confidence $|m|$ (the gap between the expected 1--9 scores); L = Qwen3-VL-32B pairwise with its own $|\lambda|$; H = a human.
Cost units: one 7--9B forward pass = 1, one 32B pass = 4, one human label = 100. P and S each cost two passes; L costs eight.
Because the 32B judge is slow, we run it on a stratified subset of at most 700 pairs per benchmark (all 1{,}860 Place Pulse pairs), and all ladder results use that subset.
Each configuration is calibrated with \cref{alg:cal} on a random half of the strict pairs and evaluated on the other half, over 50 splits, with $\delta=0.1$.
``Fixed'' ladders use the single equal-quantile path; adaptive ladders use the on/off path family of \cref{sec:method} (2 paths for P$\to$S and P$\to$L, 4 for P$\to$S$\to$L, which therefore contains the other three as special cases).
We report automation, the share of test pairs decided without a human. Mean test risk stays below $\alpha$ for every configuration and dataset (maximum .17 at $\alpha=.20$ and .26 at $\alpha=.30$).

\paragraph{Baselines.}
P alone is SCOPE \citep{badshah2026scope} with $\lambda$ in place of the equivalent entropy (\cref{sec:diagnostic}); L alone is SCOPE with the stronger judge; P$\to$L is a two-model Trust-or-Escalate cascade \citep{jung2024trust}.
We also report order-consistency filtering \citep{seo2026when}, which keeps only pairs whose verdict survives the swap. It has no risk target: it keeps 33--94\% of pairs, but the error among kept pairs is still .08--.38, for example .31 instead of .39 on MMRB2 for Qwen3-VL-8B. As a selection rule it leaves most of the error in place, and it cannot be tuned to a target.

\begin{table*}[t]
\centering\scriptsize
\setlength{\tabcolsep}{3pt}
\caption{Certified automation (\%) with Qwen3-VL-8B as the small judge (mean over 50 splits; ladder subset, $n$ = strict pairs). Bold: best per row and $\alpha$. P: pairwise; S: pointwise; L: Qwen3-VL-32B pairwise. P$\to$S (fixed) uses a single path; the others choose adaptively among rung subsets. Other small judges in \cref{app:ladders}.}\label{tab:ladder}
\begin{tabular}{l|cccccc|cccccc}
\toprule
 & \multicolumn{6}{c}{$\alpha=0.20$} & \multicolumn{6}{c}{$\alpha=0.30$}\\
Dataset ($n$) & P & P$\to$S (fixed) & P$\to$S & L & P$\to$L & P$\to$S$\to$L & P & P$\to$S (fixed) & P$\to$S & L & P$\to$L & P$\to$S$\to$L\\
\midrule
VL-RewardBench (700) & 1 & 4 & 1 & \textbf{13} & 1 & 4 & 46 & \textbf{75} & 60 & 72 & 62 & 67\\
Multimodal RB (700) & 24 & 74 & 71 & 67 & 70 & \textbf{76} & 92 & \textbf{100} & 99 & 100 & 98 & 100\\
MM-RLHF-RB (170) & \textbf{35} & 23 & 15 & 29 & 23 & 11 & 86 & 76 & 77 & \textbf{86} & 82 & 74\\
MLLM-as-a-Judge (551) & 1 & 0 & 0 & 2 & \textbf{3} & 1 & 19 & 14 & 14 & 19 & \textbf{32} & 29\\
GenAI-Arena (457) & 37 & 22 & 27 & 27 & \textbf{37} & 36 & \textbf{84} & 82 & 79 & 66 & 82 & 78\\
MJ-Bench (700) & 83 & 68 & 81 & \textbf{85} & 84 & 80 & 100 & \textbf{100} & 100 & 100 & 100 & 100\\
MMRB2 (700) & 0 & 0 & 0 & 0 & 0 & 0 & 9 & \textbf{9} & 5 & 7 & 5 & 6\\
ImageRewardDB (700) & 1 & 0 & 0 & \textbf{1} & 0 & 0 & 5 & 6 & 5 & \textbf{12} & 7 & 4\\
Q-Bench2 (442) & 84 & \textbf{92} & 92 & 30 & 76 & 92 & 100 & \textbf{100} & 100 & 56 & 100 & 100\\
EditReward-Bench (115) & 0 & 0 & 0 & 0 & 0 & 0 & 1 & \textbf{17} & 10 & 1 & 2 & 5\\
Place Pulse 2.0 (1500) & 1 & 1 & 1 & \textbf{2} & 1 & 0 & 41 & \textbf{49} & 41 & 39 & 40 & 37\\
\bottomrule
\end{tabular}

\end{table*}

\begin{table}[t]
\centering\small
\caption{Mean certified automation (\%) and mean cost per pair, over the 11 datasets. Human label = 100; ``all to humans'' would cost 100 + the cost of the first rung.}\label{tab:ladder-mean}
\begin{tabular}{llcccccc}
\toprule
Small judge & $\alpha$ & P & P$\to$S (fix) & P$\to$S & L & P$\to$L & P$\to$S$\to$L\\
\midrule
Qwen3-VL-8B & .20 & 24 & 26 & 26 & 23 & 27 & \textbf{27}\\
            & .30 & 53 & \textbf{57} & 54 & 51 & 55 & 55\\
            & cost@.30 & 49.2 & \textbf{45.9} & 49.5 & 57.3 & 50.4 & 52.3\\
Qwen3.5-9B  & .20 & 28 & 33 & 30 & 23 & 32 & \textbf{34}\\
            & .30 & 55 & \textbf{63} & 61 & 51 & 60 & \textbf{63}\\
            & cost@.30 & 47.4 & \textbf{39.5} & 42.3 & 57.3 & 45.1 & 42.9\\
Qwen2.5-VL-7B & .20 & 22 & 23 & 23 & 23 & \textbf{26} & 25\\
            & .30 & 38 & 39 & 37 & \textbf{51} & 47 & 45\\
            & cost@.30 & 63.9 & 64.7 & 66.3 & \textbf{57.3} & 59.8 & 62.8\\
\bottomrule
\end{tabular}
\end{table}

\paragraph{Result 1: no rung order dominates.}
\Cref{tab:ladder,tab:ladder-mean} show that the best configuration changes with the dataset, the judge and $\alpha$.
Averaged over datasets, the best escalating configuration adds 3--9 points of automation over pairwise-only selection and lowers mean cost (e.g.\ 47.4 $\to$ 39.5 for Qwen3.5-9B at $\alpha=.30$). The largest single gains are large: for Qwen3-VL-8B on Multimodal RB at $\alpha=.20$, automation rises from 24\% to 74--76\%; for Qwen3.5-9B on Q-Bench2 at $\alpha=.30$, from 24\% to 96\%.
The pointwise rung helps on Multimodal RB, VL-RewardBench, Q-Bench2 and EditReward-Bench, and, on the fixed path, with the Qwen3 judges on Place Pulse. It hurts on MJ-Bench and GenAI-Arena, where pointwise accuracy is 2--5 points below pairwise. The fixed P$\to$S path loses most there (MJ-Bench at $\alpha=.10$: 42\% $\to$ 8\% for Qwen3-VL-8B); the adaptive path recovers most of it (40\%).

\paragraph{Result 2: a 4$\times$ larger judge is not a better judge here.}
Qwen3-VL-32B is more accurate than the 8B model on only some datasets (MLLM-as-a-Judge .73 vs .70, Multimodal RB .80 vs .77) and less accurate on others (MMRB2 .59 vs .60, Q-Bench2 .85 vs .89); on Place Pulse it reaches .661.
As a stand-alone selective judge it is the best choice only for the weakest small judge (Qwen2.5-VL-7B, 51\% vs 38\% automation at $\alpha=.30$). For the Qwen3 judges it certifies less than the small judge's own pairwise rung, at higher cost.
As a second rung it helps where its errors differ from the small judge's, notably MLLM-as-a-Judge (18\% $\to$ 32\% for Qwen3-VL-8B at $\alpha=.30$).

\paragraph{Result 3: certification is data-hungry.}
The subset results show how strongly automation depends on calibration size. On Multimodal RB at $\alpha=.20$, pairwise-only automation for Qwen3-VL-8B is 60\% with 750 calibration pairs (full data, \cref{app:ladders}) but 24\% with 350. Because the Clopper--Pearson bound tightens slowly, a practitioner who can afford a few hundred more calibration labels may gain more than from any extra rung.

\paragraph{Result 4: the hard image-preference sets certify almost nothing.}
On ImageRewardDB no configuration exceeds 17\% automation at $\alpha=.30$; on MMRB2 only Qwen3.5-9B certifies a sizeable share (up to 49\%). At $\alpha=.20$ no configuration reaches 9\% on either. The base error of these judges (35--44\%) is too close to $\alpha$, and confidence ranks errors too weakly (AUROC .57--.65) to find a low-risk subset.
For these tasks, today's open VLM judges are suitable only inside estimators that do not need per-pair trust, such as the PPI rankings of \cref{sec:ranking}.
